\title{Large Language Models Can Automatically Engineer Features for Few-Shot Tabular Learning}

\begin{abstract}
Large Language Models (LLMs) have demonstrated significant capability in addressing complex and novel reasoning challenges, suggesting their strong suitability for tabular data tasks, which are critical across numerous practical domains. In this work, we introduce an innovative in-context learning paradigm named \textsf{FeatLLM}, designed to leverage LLMs as automatic feature constructors that transform the original input into a representation optimally tailored for tabular prediction tasks. These generated feature sets are subsequently utilized by a straightforward downstream model, such as linear regression, to accurately estimate class probabilities, achieving robust performance even in few-shot scenarios. Notably, \textsf{FeatLLM} restricts prediction-time computation to the simple model acting on the learned features, dispensing with the requirement to interrogate the LLM per sample during inference. Additionally, our method is compatible with basic API-level LLM access and sidesteps issues of prompt length constraints. Through empirical analysis across various tabular datasets spanning different fields, \textsf{FeatLLM} is shown to craft effective rule-based features and achieves considerable improvements—typically exceeding 10\% in accuracy—compared to competitors like TabLLM and STUNT.
\end{abstract}

\section{Introduction}
Recent advances have positioned Large Language Models (LLMs) as powerful tools for handling novel tasks and delivering high-quality responses even in the absence of task-specific training~\cite{creswell2022selection,imani2023mathprompter,taylor2022galactica}. Leveraging large-scale text corpora, LLMs inherently capture deep and broad prior knowledge, allowing them to fill the role traditionally occupied by subject-matter experts across varied settings~\cite{Genomes2021,singhal2023large,wilcox2003role}. This versatility has fueled their adoption for applications such as conversational analysis~\cite{finch2023leveraging} and automated program correction~\cite{fan2023automated}. Moreover, by incorporating explicit task instructions and illustrative examples into their prompts, LLMs are able to contextualize their understanding, focusing on the relevant features and instances for the task at hand~\cite{brown2020language,zhao2023survey}. Such abilities highlight their promise in automating feature engineering through contextual data interpretation.

Building upon these strengths, researchers have begun to apply the domain knowledge and inferential power of LLMs to tabular learning challenges. For example, incorporating background knowledge—such as recognizing the association between age and disease risk—can enhance disease prediction models. Contemporary approaches encode natural language descriptions of both tasks and features, serialize the structured input, and then process this through an LLM~\cite{dinh2022lift,wang2023anypredict}. These strategies frequently make use of in-context learning~\cite{nam2023semi} or lightweight adaptation methods like parameter-efficient fine-tuning~\cite{dinh2022lift,hegselmann2023tabllm}. Empirical studies have found that the knowledge embedded in LLMs is particularly advantageous in low-data scenarios, often surpassing the performance of traditional baselines for tabular prediction~\cite{hegselmann2023tabllm}.

Existing LLM-based approaches to tabular data analysis exhibit several notable shortcomings. In end-to-end prediction pipelines, each data instance necessitates at least one forward pass through the LLM, resulting in significant computational overhead. Furthermore, attaining high predictive accuracy frequently hinges on fine-tuning the LLM—a strategy that becomes impractical when working with LLMs that lack publicly available training implementations or accessible fine-tuning APIs. This is particularly relevant since many recently developed high-performing LLMs only support restricted usage through inference APIs~\cite{anil2023palm,openai2023gpt4}. Additionally, the efficacy of most current methods diminishes with longer prompts~\cite{liu2023lost}; as tabular datasets grow in the number of features, the resulting serialized inputs to the LLM can easily exceed prompt limitations or suffer in terms of performance and efficiency. Such limitations substantially hinder the deployment of LLM-centric tabular models in practical scenarios.

To address these issues, we advance a novel strategy that reframes the function of LLMs from direct prediction to sophisticated feature construction. Rather than employing LLMs exclusively for end-to-end inference, we utilize their capacity for integrating prior knowledge by delegating feature engineering tasks to them. In particular, we propose \textsf{FeatLLM}, an innovative in-context learning paradigm focused on elucidating the underlying ``criteria" that inform LLM predictions. For example, in a medical diagnosis context, the LLM can articulate interpretable rules specifying which conditions on feature values signal the presence of a specific disease. These derived rules facilitate the synthesis of new, task-relevant features to replace original attributes, thus simplifying subsequent modeling. By offloading complex decision-making to the feature construction phase, downstream learning is considerably streamlined, leading to faster inference as resource-intensive modeling becomes unnecessary. Notably, this framework capitalizes on LLMs' in-context learning prowess—features are extracted solely via illustrative prompt demonstrations, obviating the requirement for explicit model training. Moreover, by incorporating ensemble-inspired techniques such as feature bagging~\cite{breiman2001random}, our approach effectively contends with the prompt length constraints encountered in high-dimensional tabular settings.

\textsf{FeatLLM} divides the process of feature engineering via prompt construction into two distinct stages: (1) interpreting the core problem while deducing how features relate to the prediction targets; and (2) using insights from the prior analysis, along with a limited set of examples, to formulate actionable rules that distinguish among classes. This stepwise approach encourages the LLMs to disregard irrelevant or redundant features during rule formulation. Once inferred, these rules are mapped to the dataset (by translating them into executable code), transforming each instance into binary indicators reflecting compliance with each rule. A simple, interpretable model is subsequently trained on these derived features to predict class probabilities. Additionally, the method’s robustness is enhanced through ensembling, involving the repeated creation of features while incorporating feature bagging strategies. By carefully selecting how many features and samples participate in each bag, the framework addresses constraints related to large prompt sizes. Free from the limitations of retraining and maintaining a low computational footprint at inference time, \textsf{FeatLLM} makes large language model-driven feature engineering applicable across a spectrum of real-world applications.

Our experiments with \textsf{FeatLLM} span 13 tabular datasets under low-resource scenarios, demonstrating its consistent and superior performance. The results reveal the LLMs' effective integration of background and in-dataset knowledge in extracting discriminative features. Compared to recent few-shot learning alternatives, our framework achieves stronger results in varied configurations. The implementation is available through an anonymized GitHub repository at~\texttt{https://github.com/Sungwon-Han/FeatLLM}.

\section{Related Work}

\subsection{Few-Shot Learning with Tabular Data}

Progress in few-shot learning has made it possible to extract transferable representations from datasets with only limited annotated examples, reducing labeling demands~\cite{chen2018closer}. Although few-shot learning originated in the context of image analysis, its applicability has now been established across a spectrum of data types, such as text, audio, and, more recently, tabular formats~\cite{majumder2022few,nam2023stunt,schick2021exploiting}. A persistent issue when relying on only a handful of labeled instances is the propensity to pick up on incidental correlations~\cite{bartlett2021deep}. To mitigate this, researchers have introduced auxiliary information and domain-centric prior knowledge to help embed relevant inductive biases during model training. Notably, some strategies supplement scarce labeled data with large unlabeled collections for semi-supervised training, leveraging careful augmentation~\cite{ucar2021subtab,yoon2020vime}. Others target improved model starting points via unsupervised pre-training, disconnected from specific downstream tasks~\cite{somepalli2021saint}. Unsupervised learning methods frequently use objectives based on reconstructing data after masking or corruption~\cite{arik2021tabnet,majmundar2022met} or apply data augmentation to facilitate contrastive learning~\cite{bahri2022scarf,chen2020simple}. However, because tabular data tends to be highly heterogeneous, crafting universally effective data augmentations proves difficult, which has, in turn, limited the impact of these approaches under few-shot scenarios~\cite{nam2023stunt}.

More recently, methods have been introduced where models are pre-trained across a variety of tabular datasets, not only those directly linked to the target domain, before employing transfer learning on the intended task~\cite{zhang2023generative,levin2022transfer,zhu2023xtab}. An illustrative example, TransTab~\cite{wang2022transtab}, utilizes transformer architectures that learn semantic associations not solely from table cell values, but also via column names, helping to transfer information between otherwise unrelated tabular sources. This technique facilitates both zero-shot and few-shot adaptation to novel tabular learning problems. In a different vein, TabPFN~\cite{hollmann2022tabpfn} constructs synthetic tabular examples from mixed data distributions for the unsupervised pre-training of a Bayesian neural network. At deployment, this network can process both the labeled and test data together without further optimization. This acquired general knowledge from diverse data sources consistently leads to superior results compared to traditional decision-tree approaches, demonstrating efficacy in few-shot learning applications.

\subsection{Language-Interfaced Tabular Learning}

Incorporating large language models (LLMs)~\cite{anil2023palm,openai2023gpt4} represents a promising avenue for utilizing external knowledge in tabular data tasks. With their foundation in massive, varied text corpora, LLMs have shown a robust capacity to generalize, successfully handling numerous and previously unseen tasks~\cite{brown2020language}. Recent investigations have sought to leverage LLMs’ extensive encoded priors for applications in tabular data analysis. A common approach involves converting tabular data into a text format, then supplying this serialized input, potentially accompanied by relevant metadata and task details, to an LLM via prompts~\cite{dinh2022lift,hegselmann2023tabllm,wang2023anypredict}. By presenting table content along with feature and task information in textual form, LLMs can infer key task-feature interactions and perform inference accordingly. Developments in this line of work have included training LLMs using efficient adaptation methods such as LoRA~\cite{hu2021lora} or IA3~\cite{liu2022few}, and exploring prompt-based in-context learning, which enables LLMs to leverage a handful of demonstration examples embedded in the prompt—eliminating the need for additional training~\cite{dinh2022lift,hegselmann2023tabllm,nam2023semi,slack2023tablet}.

Although prior methods have shown promise in enhancing few-shot tabular learning, their reliance on routing inference through an LLM at all times introduces practical constraints for deployment in industrial scenarios. This research seeks to diverge from the traditional black-box model that computes predictions for every instance exclusively via an LLM. Instead, it directs the LLM to deduce the distinguishing rules or criteria for each class. \textsf{FeatLLM} takes the rules generated by the LLM, converts them into executable program code, and generates new features for the input data—thus enabling predictions to be made without subsequent calls to the LLM. This harnesses the power of in-context learning to derive rules by utilizing existing knowledge and the supplied few-shot samples.

\section{Methods}
\textsf{FeatLLM} focuses on mining ``rules," or characteristic conditions for each potential output class, based on a limited number of supplied examples, as illustrated in Figure~\ref{fig:main_model}. These rules, once articulated by the LLM, are parsed and translated into binary feature indicators corresponding to whether each rule holds true for a given data point. The resulting collection of binary features is multiplied by a set of non-negative trainable weights through a dot product to estimate class probabilities. By optimizing these weights during training, the model discerns which rules exert the greatest influence on classification outcomes (Section~\ref{Sec:training}). This workflow is iterated several times using bagging, with predictions subsequently aggregated through ensembling. The sections that follow articulate the problem statement and elaborate on the specifics of each stage.

\vspace*{-0.12in}
\paragraph{Problem formulation.} Consider a tabular dataset consisting of $N$ labeled examples, represented as $\mathcal{D} = \{(\mathbf{x}^i, \mathbf{y}^i)\}_{i=1}^{N}$. Each entry $\mathbf{x}^i$ corresponds to a $d$-dimensional feature vector, where the set of feature names $F = \{f_j\}_{j=1}^{d}$—given in natural language—are provided, along with separate definitions for these features. In a classification scenario, the label $\mathbf{y}^i$ belongs to a predefined list of classes. For $k$-shot learning scenarios, the model is trained only on $k$ labeled instances randomly drawn from the total $N$.

\subsection{Prompt Design for Extracting Rules}
\label{Sec:prompt_design}
To facilitate rule extraction from LLMs using a reasoning path reflective of domain expertise, we structure the prompt to encourage the model to approach tabular learning problems in an expert-like manner. The prompt is composed of three principal elements (see Fig.~\ref{fig:prompt_for_rule} and Appendix~\ref{sec:example_rule}):

\vspace*{-0.12in}
\paragraph{Basic information description.} This segment conveys the core data required for problem resolution (see the orange-highlighted section in Fig.~\ref{fig:prompt_for_rule}), including the specification of the task accompanied by label information, the explanations of feature variables, and sample instances with annotations. The task prompt takes the form of a question (e.g., ``Does this patient's myocardial infarction complications data indicate chronic heart failure? Yes or no?"). Additional sample prompts are included in Appendix~\ref{sec:task_desc}. Each feature is characterized by its type (e.g., categorical, continuous), and a definition may be optionally attached. For categorical features, a subset of example categories may be shown. Representative examples from the training set are serialized as textual descriptions together with their correct labels. Serialization adheres to protocols established in earlier works~\cite{dinh2022lift,hegselmann2023tabllm}:

\begin{align}
&\text{Serialize}(\mathbf{x}^i,\mathbf{y}^i, F) = \nonumber \\ 
&``\ f_1\ \text{is}\ \mathbf{x}^i_1.\ ... \ f_d\  \text{is}\  \mathbf{x}^i_d.\ \ \text{Answer:}\  \mathbf{y}^i\ ”
\end{align}

\vspace*{-0.12in}
\paragraph{Reasoning instruction.} Our objective is to strengthen the LLM's reasoning capabilities by supplying explicit reasoning instructions (refer to the blue segments in Fig.~\ref{fig:prompt_for_rule}). The instructions initiate with an introductory statement akin to the chain-of-thought paradigm~\cite{wei2022chain}. Subsequently, we delineate the reasoning process into two stages. In the initial stage, the LLM is prompted to deduce potential causal linkages or tendencies between the input features and the overall task specification, independent of example demonstrations and instead grounded in general background knowledge or intuition. This facilitates comprehensive utilization of the LLM’s prior understanding of the task. In the subsequent stage, the LLM integrates knowledge from sample demonstrations and the reasoning generated in the first stage to formulate rules pertinent to each class. Employing this bifurcated process mitigates the risk of the model learning spurious associations from irrelevant variables and sharpens its attention on features with substantive relevance. To avoid scenarios of rule sparsity or excess—which could result in underfitting or unnecessarily large prompts—we fix the rule count for each class to a constant value (specifically, 10)\footnote{A discussion on optimal rule number is available in Section~\ref{sec:analysis}.}. Furthermore, during the rule formulation process, the LLM is instructed to consider feature types as specified in the accompanying basic information segment to preclude inconsistencies due to type errors.

\vspace*{-0.12in}
\paragraph{Response instruction.} To streamline the subsequent interpretation and application of the deduced rules, we employ targeted response instructions to specify the expected structure of the LLM’s output (see the yellow sections in Fig.~\ref{fig:prompt_for_rule}). These prompts clarify both the general template for organizing responses for each answer class (i.e., expected response format) and the proper syntax that individual rules should adhere to (i.e., expected rule format). By offering this guidance, the LLM can adjust the format according to the specific feature types in question. In the absence of restrictive guidance, the LLM is permitted to compose compound rules by employing logical connectors such as AND or OR. Illustrative examples of expected outputs are detailed in Appendix~\ref{sec:outcome-rule}.

\subsection{Inferring Class Likelihood via Rules}
\label{Sec:training}

\paragraph{Parsing rules for feature generation.} In this stage, we utilize the previously generated rules to derive new binary features for each sample. For every class, a binary indicator is assigned to reflect whether or not the sample adheres to that class's rules—coded as '0' or '1'. These features serve as proxies for the class membership likelihood of each sample. Since the LLM-generated rules are in natural language, a systematic parsing approach is necessary for feature extraction. Although response and rule formatting are standardized during generation to facilitate easier parsing, outputs from the LLM may still be affected by inconsistent syntactic elements~\cite{kovavc2023large}. For example, the LLM might express inequalities verbally (e.g., replacing ``$>$” or ``$<$” with ``greater than'' or ``less than''), which introduces further complexity into the parsing process.

To overcome the difficulties presented by noisy textual output, we refrain from developing elaborate custom code. Instead, we exploit the LLM's proficiency in interpreting semantic intent across varying syntactic forms, alongside its capability to produce code according to given prompts (thanks to its training on code-related data). Accordingly, the input prompt to the LLM includes the function’s name, descriptions of inputs and outputs, as well as the rules to be implemented. This input is fed directly into the LLM. Illustrative prompts and their corresponding code outcomes are provided in Figures~\ref{fig:prompt_for_parsing} and~\ref{sec:example_parsing}-\ref{sec:example_parse_outcome} in the Appendix. The resulting code is then executed using Python’s \texttt{exec()} functionality for dataset transformation.

\vspace*{-0.12in}
\paragraph{Inferring class likelihood.} With binary features derived from rules specified for each class, one straightforward approach to quantify a sample’s likelihood of class membership is tallying the count of satisfied rules for every class (equivalent to summing binary features per class). However, the predictive strength of individual rules and features may vary, necessitating that their relative contributions be estimated from the training data. We approach this by training a conventional machine learning (ML) algorithm on the binary features for each class, thereby learning their respective importances. As an instance, consider a bias-free linear model. Let the binary feature vector for sample $\mathbf{x}^i$ corresponding to class $k$ be denoted by $\mathbf{z}_k^i \in \{0, 1\}^R$, where $R$ indicates the number of rules for that class, and $c$ the total number of classes. The probability vector for each class $\mathbf{p}^i$ can be evaluated as:

\begin{align}
\text{logit}_k^i &= \max(\mathbf{w}_k, 0) \cdot \mathbf{z}_k^i, \nonumber \\
\mathbf{p}^i &= \text{Softmax}({[\text{logit}_{1}^i, ..., \text{logit}_{c}^i]}). \label{eq:infer_prob}
\end{align}

In this approach, $\mathbf{w}_k$ denotes the tunable parameter in the linear model, which encodes the relative influence of each feature. By constraining these weights within a positive subspace, we guarantee that features sourced from the LLM do not detract from a class’s probability during model training. The vector of logits is subsequently passed through the softmax transformation to produce class probabilities. The weight vector $\mathbf{w}_k$ is optimized by minimizing cross-entropy loss using a limited set of training instances.

\vspace*{-0.12in}
\paragraph{Ensembling with bagging.} To produce the final classification, we replicate the entire inference pipeline multiple times, yielding several independent models whose outputs are then ensembled. The learned parameters from $T$ runs, $\{\mathbf{w}[t]\}_{t=1}^T$, are leveraged by computing each class probability $\mathbf{p}^i$ using the corresponding $\mathbf{w}[t]$ for each run, followed by averaging these probabilities as prescribed in Eq.~\ref{eq:infer_prob}.

Randomness is infused into the rule generation phase by setting a sufficiently elevated temperature during LLM inference, promoting diverse rule extraction. Additionally, we permute the arrangement of few-shot examples, motivated by findings that their order can alter LLM output~\cite{lu2022fantastically}\footnote{Further exploration of rule variability appears in Appendix~\ref{sec:diversity}}. To further exploit the diversity induced in predictions and improve stability, we implement bagging, selecting random subsets of features or instances for each trial—a strategy particularly beneficial with more training examples or higher-dimensional input. The suggested ensembling method offers two notable strengths: first, any errors from erroneous LLM-derived rules in a subset of trials are mitigated by correct rules generated in other runs, thus bolstering framework robustness. This effect is associated with LLM self-consistency~\cite{wang2022self}, as rules repeatedly produced across runs tend to be reliable. Second, ensemble bagging alleviates constraints on LLM prompt length; if the size of the tabular dataset exceeds the LLM’s input capacity, bagging can trim input size per trial while maintaining robust performance. While any individual rule set may omit information about some features, the aggregate effect of combining numerous weak learners—each drawn from different feature or instance samples—minimizes shortcomings due to partial coverage in single trials.

\section{Experiments}
We assess the effectiveness of \textsf{FeatLLM} on a diverse range of tabular datasets, alongside an array of analytical studies to gain insight into the mechanisms underpinning our approach. Details of further experimental results—including tests using alternative LLMs, qualitative studies, and others—are provided in the Appendix, owing to limitations of space.

\subsection{Performance Evaluation}
\paragraph{Datasets.}
Our study makes use of 13 datasets intended for binary or multi-class classification problems: (1) Adult~\cite{asuncion2007uci}, which involves determining whether a person's annual income exceeds \$50,000; (2) Bank~\cite{moro2014data}, focused on forecasting whether clients will purchase a term deposit; (3) Blood~\cite{yeh2009knowledge}, used to classify if donors will make subsequent donations; (4) Car~\cite{kadra2021well}, for assessing vehicle quality; (5) Communities~\cite{misc_communities_and_crime_183}, centered on estimating regional crime rates; (6) Credit-g~\cite{kadra2021well}, which predicts a person's credit risk category; (7) Diabetes\footnote{\texttt{kaggle.com/datasets/uciml/pima-indians-diabetes-database}}, designed to detect the presence of diabetes; (8) Heart\footnote{\texttt{kaggle.com/datasets/fedesoriano/heart-failure-prediction}}, for forecasting the likelihood of coronary artery disease; and (9) Myocardial~\cite{misc_myocardial_infarction_complications_579}, concerning prediction of chronic heart failure.

In addition to these, we include both recently-released and synthetic datasets, which are infrequently explored in prior work and were absent from the LLM pre-training corpora: (10) Cultivars\footnote{\texttt{archive.ics.uci.edu/dataset/913}}, which involves predicting whether the grain yield of a particular soybean cultivar is high; (11) NHANES\footnote{\texttt{archive.ics.uci.edu/dataset/887}}, tasked with classifying a subject’s age group; (12) Sequence-type, concerned with categorizing a sequence as either arithmetic, geometric, Fibonacci, or Collatz; and (13) Solution-mix, where the objective is to predict if a mixture’s concentration exceeds 0.5. The first two of these datasets were newly added to the UCI Machine Learning Repository after September 2023, ensuring that the LLM API (gpt-3.5-turbo-0613) applied in our framework could not have included them in pre-training. The last two datasets are artificially generated, further eliminating the risk of prior memorization by the LLM and guaranteeing unbiased evaluation.

The included datasets demonstrate considerable heterogeneity in both scale and difficulty, as summarized in Table~\ref{Tab:basic-desc}. Each resource is accompanied by explicit names and attribute descriptions. Notably, datasets such as `Communities' and `Myocardial' feature in excess of 100 variables, which is particularly challenging for LLMs given the constraints of limited context window sizes.

\vspace*{-0.12in}
\paragraph{Baselines.}
For our comparative analysis, we benchmark against nine reference methods. The initial three represent traditional models for tabular data: (1) Logistic regression (LogReg), (2) XGBoost~\cite{chen2016xgboost}, and (3) RandomForest~\cite{ho1995random}. The next two—(4) SCARF~\cite{bahri2022scarf} and (5) STUNT~\cite{nam2023stunt}—leverage unlabeled data, although in actual practice, obtaining extensive unlabeled collections can be challenging. The sixth method, (6) TabPFN~\cite{hollmann2022tabpfn}, employs synthetic data generated across a variety of distributions as pre-training material. Finally, the last three make use of LLMs: (7) In-context learning (In-context)~\cite{wei2022emergent}, (8) TABLET~\cite{slack2023tablet}, and (9) TabLLM~\cite{hegselmann2023tabllm}. In-context learning works by injecting a handful of training exemplars into the prompt directly without explicit model tuning. TABLET enhances this by supplementing the prompt with auxiliary signals obtained from external rule-based classifiers. TabLLM adopts a parameter-efficient tuning strategy, such as IA3~\cite{liu2022few}, for few-shot fine-tuning of the LLM.

\vspace*{-0.12in}
\paragraph{Implementation details.}
\textsf{FeatLLM} is built atop GPT-3.5 as its primary LLM, though the framework itself is compatible with various LLM architectures (see Appendix~\ref{sec:backbones} for comparative analyses using alternative backbones). During LLM inference, the temperature parameter is fixed at 0.5, and the top-p parameter retains the API’s default setting of 1. The ensemble count is specified as 20, with 10 extraction rules employed. Additional insights regarding the influence of hyper-parameters can be found in Figure~\ref{fig:hyperparameter2} and Appendix~\ref{sec:hyper-parameter}. For the linear model, we optimize using the Adam optimizer at a learning rate of 0.01 across 200 training epochs, employing $k$-fold cross-validation to determine the best epoch. When implementing baselines, in-context learning leverages GPT-3.5, whereas TabLLM uses T0~\cite{sanh2022multitask}, as prescribed in its original configuration. It is important to note that TabLLM confines the choice of LLMs to those with publicly released checkpoints, thereby excluding GPT-3.5. For conventional machine learning algorithms (LogReg, XGBoost, RandomForest), hyper-parameters are tuned by Grid-search in conjunction with $k$-fold cross-validation. The $k$ value is chosen as 2 or 4, ensuring that each class appears at least once in every training split. Further implementation specifics are detailed in Appendix~\ref{sec:baselines}.

\vspace*{-0.12in}
\paragraph{Results.}
Tables~\ref{tab:main1} and \ref{tab:main2} report a comprehensive comparison of model performances across diverse datasets. Each experiment is independently executed three times with varied random splits for training and testing sets; we then calculate the mean and standard deviation of AUC scores. Across nearly all settings, \textsf{FeatLLM} either outperforms all baselines or achieves the second-best results. Remarkably, our method with only 4-shot scenarios can match the accuracy of traditional tabular models that use 32 samples, as illustrated in Fig.~\ref{fig:compares}. While STUNT and TabPFN deliver notable gains on particular datasets, their aggregate advantage over standard machine learning methods remains limited. Furthermore, other LLM-based models—such as in-context learning, TABLET, and TabLLM—struggle to produce results on tabular data containing a large feature set. In contrast, our approach effectively handles datasets with many features via combined bagging and ensemble strategies, delivering robust and reliable outcomes. It is worth mentioning that our ensemble methodology could, in theory, be incorporated into other LLM-driven methods. However, doing so would double the number of inference operations per sample, significantly increasing computational burden, and consequently rendering such approaches less practical.

\subsection{Analysis \& Discussion}
\label{sec:analysis}
\paragraph{Ablation study.} To assess how individual framework components affect overall performance, we conduct ablation studies, examining: (1) \textsf{-Tuning}: removing the linear model's weight adjustment stage, instead aggregating binary feature outputs to calculate the logit; (2) \textsf{-Ensemble}: excluding the ensemble method; (3) \textsf{-Description}: removing the use of feature descriptions; and (4) \textsf{-Reasoning}: omitting Step-1 in the reasoning instruction phase during rule creation.

Table~\ref{tab:ablation} outlines the average AUC changes resulting from these ablations across all datasets. Disabling any individual component consistently causes a decrease in predictive performance. Notably, weight tuning has an increasingly positive effect as the number of shots grows, suggesting that abundant data allows for better estimation of rule relevance, making tuning more beneficial. Conversely, for low-shot scenarios, the contributions from feature descriptions and reasoning steps are magnified, indicating that leveraging LLM’s inherent prior knowledge is especially valuable in data-scarce regimes. Importantly, the ensemble mechanism proposed here consistently yields improved performance regardless of the configuration.

\vspace*{-0.12in}
\paragraph{Training \& inference time comparison.} We analyze and compare both training and inference durations across different approaches. All tests utilize the Adult dataset, with the training data comprising 16 examples. A single A100 GPU serves as the standard computation resource, with the exception of TabLLM, which leverages four A100 GPUs in a model parallelism configuration. Table~\ref{tab:cost} outlines the timing statistics. Of particular note, \textsf{FeatLLM} achieves inference times that are relatively low and similar to those of traditional baselines. Conversely, alternative LLM-oriented techniques—including In-context learning, TABLET, and TabLLM—are characterized by significantly longer inference times. For training times, \textsf{FeatLLM} exhibits performance on par with other LLM-backed methods, requiring just 30 API invocations. This number is minimal and not cost-prohibitive and could further be reduced through parallelization.

\vspace*{-0.12in}
\paragraph{Quality of features generated by \textsf{FeatLLM}.}
To assess how informative the rule-based features produced by \textsf{FeatLLM} are in real-world settings, we ran a straightforward experiment. Here, we calculate the mean AUROC of the generated features with respect to the target labels for each dataset, and report the proportion of features with an AUROC above 0.5 as indicative of utility. The summarized findings are presented in Table~\ref{tab:quality}. These results indicate that most constructed rules align well with the target variables, yielding substantial utility for the problem at hand (see the ``Before tuning'' column). Furthermore, as part of the linear model fitting phase, rules deemed irrelevant—those assigned weights below a certain threshold—are systematically excluded (see the ``After tuning'' column). For these experiments, a weight threshold of 0.1 was chosen, based on inspection of the weight distribution.

\vspace*{-0.12in}
\paragraph{Effect of adding spurious correlations.} To evaluate how well different approaches remove extraneous spurious correlations when limited labeled data is available, we performed an empirical assessment. Specifically, we focused on the Heart dataset, randomly augmenting it with columns drawn from the Adult dataset that are not pertinent to the Heart dataset’s prediction task. Subsequently, we examined how model performance evolves as the count of these irrelevant columns grows. To maintain experimental fairness, all methods are provided the same 8 labeled examples and do not utilize any unlabeled data, which means methods like SCARF and STUNT are omitted from this analysis.

Figure~\ref{fig:spurious} presents the observed model performances as more spurious features are introduced. Among all tested models, \textsf{FeatLLM} exhibited the smallest decrease in performance in response to increased spurious correlations. This resilience is presumably attributable to our framework’s design, which leverages domain knowledge to assess the alignment between each feature and the primary task during rule extraction, thereby avoiding the inclusion of irrelevant features in generated rules and leading to greater robustness. Empirically, rules influenced by spurious columns made up only 1-2\% of the rule set. On the other hand, it is noteworthy that \textsf{FeatLLM} may inadvertently capture spurious associations and misinterpret the underlying causal links between target and features if irrelevant columns are sourced from datasets that are themselves closely related to the original task (see Appendix~\ref{sec:spurious-appendix}).

\vspace*{-0.12in}
\paragraph{Hyper-parameter analysis}
\textsf{FeatLLM} relies primarily on two hyper-parameters: the ensemble size and the number of rules generated per class in each LLM query. We systematically examined how variations in these settings influence overall model effectiveness. Our findings indicate that increasing the ensemble count yields better outcomes, yet returns diminish past a threshold (approximately 20), after which further increases offer limited incremental benefit (refer to Fig.~\ref{fig:appendix-hyperparameter1} in Appendix). In terms of rule count, although our analysis did not reveal statistically significant effects, selecting 10 rules achieved an optimal trade-off between prompt complexity and predictive accuracy (see Fig.~\ref{fig:hyperparameter2}). We hypothesize that incorporating additional rules expands the feature space and augments available information, but in data-constrained (low-shot) scenarios, this expansion can also induce overfitting.

\section{Conclusion}
We introduce \textsf{FeatLLM}, a method that substantially enhances LLM-based approaches for few-shot tabular data analysis. Unlike approaches that treat LLMs merely as prediction engines for individual records, \textsf{FeatLLM} leverages LLMs for formulating predictive rules and performing feature engineering. Through a process of rule extraction combined with ensemble bagging, our framework minimizes inference overhead, operates solely with API calls (obviating further model training), and effectively manages data with variable numbers of features. The outcome is robust predictive performance, positioning \textsf{FeatLLM} as a strong, data-efficient option for deploying LLMs on practical tabular tasks.

\textsf{FeatLLM} is tailored to excel in low-shot learning contexts. Looking ahead, we aim to broaden its capacity for constructing advanced features applicable to larger datasets, adapt it to manipulate feature types beyond rule-based representations, and enhance its interpretability—steps that will allow broader applicability across diverse data-centric applications.

\section{Broader Impact}
Our work, \textsf{FeatLLM}, proposes to harness the extensive prior knowledge embedded within LLMs to surpass the limits of traditional supervised tabular learning—especially where data is scarce. By drawing upon pretrained knowledge, our approach can address overfitting challenges commonly faced in data-limited scenarios. This framework holds particular promise for practitioners in sectors like finance and healthcare, where expert-labeled data is both costly and difficult to obtain, yet valuable domain-specific information is abundant in pretrained LLMs. Nonetheless, as a consideration for widespread implementation, it is essential to evaluate how societal biases and inaccuracies present in LLMs’ priors might impact model decisions—since, in some cases, this background knowledge may substantially shape or even override insights drawn from observed labeled samples.

\end{document}